% !TEX root =./ECE448_Practice_main.tex
%
% L12 practice set (CW). Every question is a step in the L12 link: Morse timing
% and words per minute, decoding marks and spaces, key clicks and edge time,
% the product detector and the CW filter, and setting the threshold.
% Blank copy: \begin{solution}[h] reserves h of work space.
% SOLUTIONS copy (\iskey=1) prints the red worked solutions.

\fullwidth{\textbf{Learning Objective 6.1: } I can transmit and receive CW (Morse) using
GNU Radio and decode a received CW stream.
\\[4pt]\normalfont\small Show your work. Write a \textbf{Documentation} line at the
end (\textbf{None} if you did not collaborate).}

\begin{questions}

%%%%%%%%%%%%%%%%%%%%%%%%%%%%%%%%%%%%%%%%%
\question \textbf{Morse timing.} A dit lasts one unit of $1.2/W$ seconds at $W$ words
per minute. A dah lasts 3 units, and the key stays up for 1 unit inside a letter, 3
between letters, and 7 between words. SOS is \texttt{... --- ...}.
\begin{parts}

	\begin{minipage}[t]{0.58\linewidth}
		\part Find the length of a dit, a dah, and the gap between words at 15 WPM.
		\begin{solution}[1.0in]
		\eq{ t_\text{dit} &= \dfrac{1.2\ \text{s}}{15} = 80\ \text{ms} \\
		t_\text{dah} &= 3 \times 80 = 240\ \text{ms} \\
		t_\text{word} &= 7 \times 80 = 560\ \text{ms} }
		\end{solution}
	\end{minipage}\hfill%
	\begin{minipage}[t]{0.37\linewidth}
		\raggedleft \vspace{12pt}
		dit: \ansbox[1.0in]{80 ms}

		\vspace{4pt}
		dah: \ansbox[1.0in]{240 ms}

		\vspace{4pt}
		word gap: \ansbox[1.0in]{560 ms}
	\end{minipage}

	\begin{minipage}[t]{0.58\linewidth}
		\part Count the units in SOS, including the 7-unit gap after it, and find how
		long it takes to send at 15 WPM.
		\begin{solution}[1.2in]
		Each S is three dits with two 1-unit gaps, 5 units. O is three dahs with two
		gaps, 11 units. Two letter gaps add 6 and the word gap 7:
		\eq{ 5 + 11 + 5 + 6 + 7 &= 34\ \text{units} \\
		34 \times 80\ \text{ms} &= 2.72\ \text{s} }
		\end{solution}
	\end{minipage}\hfill%
	\begin{minipage}[t]{0.37\linewidth}
		\raggedleft \vspace{12pt}
		units: \ansbox[1.0in]{34}

		\vspace{4pt}
		time: \ansbox[1.0in]{2.72 s}
	\end{minipage}

	\part A beacon sends its three-letter identifier at 7 WPM. In words, explain why
	the standard word PARIS lets us state its speed without knowing which letters it
	sends.
		\begin{solution}[1.0in]
		WPM defines the unit, not the letters: one WPM sends the 50 units of PARIS once a
		minute, so at 7 WPM every unit lasts $1.2/7 = 171$ ms whatever the letters are.
		The identifier takes a different number of units than PARIS, so it does not
		take exactly a seventh of a minute, but its dits, dahs, and gaps all follow
		from that one unit.
		\end{solution}

\end{parts}

\newpage
%%%%%%%%%%%%%%%%%%%%%%%%%%%%%%%%%%%%%%%%%
\question \textbf{Decoding marks and spaces.} The L12 decoder is set to 15 WPM, so one
unit is 80 ms. It calls a mark under 2 units a dit and a longer one a dah, ends a letter
when a space reaches 2 units, ends a word when it reaches 5, and ignores any change
that lasts less than 0.3 unit. A human operator sends the following, in
milliseconds, starting with a mark:

\begin{center}
\renewcommand{\arraystretch}{1.2}
\begin{tabular}{| l | c | c | c | c | c | c | c | c | c | c | c |}
\hline
Mark & 230 & & 255 & & 245 & & 240 & & 75 & & 250 \\\hline
Space & & 85 & & 70 & & 260 & & 95 & & 80 & \\\hline
\end{tabular}
\end{center}

and then the key stays up for 700 ms.

\begin{parts}

	\begin{minipage}[t]{0.58\linewidth}
		\part Find the decoder's boundaries in milliseconds, and decode the text.
		\begin{solution}[1.5in]
		A length of $2 \times 80 = 160$ ms separates a dit from a dah, and a gap inside a
		letter from a gap between letters; $5 \times 80 = 400$ ms ends a word, and $0.3 \times 80 = 24$ ms is the
		noise limit. The marks of 230, 255, and 245 ms are dahs, and the 260 ms space
		ends the letter: \texttt{---} is O. Then 240, 75, and 250 ms read \texttt{-.-},
		which is K, and the 700 ms space ends the word. The text is OK.
		\end{solution}
	\end{minipage}\hfill%
	\begin{minipage}[t]{0.37\linewidth}
		\raggedleft \vspace{12pt}
		\ansbox[1.1in]{OK}
	\end{minipage}

	\begin{minipage}[t]{0.58\linewidth}
		\part The operator pauses, and the 260 ms space stretches to 420 ms. What does
		the decoder print now?
		\begin{solution}[0.8in]
		420 ms is more than 400 ms, so the decoder ends a word there and prints
		\texttt{O K} with a space between the letters.
		\end{solution}
	\end{minipage}\hfill%
	\begin{minipage}[t]{0.37\linewidth}
		\raggedleft \vspace{12pt}
		\ansbox[1.1in]{O K}
	\end{minipage}

	\part During the 700 ms space, a noise burst holds the key down for 15 ms. Explain
	what the decoder does with it, and what it would do with a 60 ms burst.
		\begin{solution}[1.0in]
		15 ms is less than the 24 ms noise limit, so the decoder ignores it and the space
		continues. A 60 ms burst lasts longer than 24 ms, so the decoder accepts it as
		a mark under 160 ms, a dit, and if the key then stays up long enough, it prints
		an extra E.
		\end{solution}

\end{parts}

\newpage
%%%%%%%%%%%%%%%%%%%%%%%%%%%%%%%%%%%%%%%%%
\question \textbf{Key clicks.} Lesson 12 measured the keyer's spectrum at 20 WPM. With
raised-cosine edges, the level 1 kHz from the carrier, compared with hard keying, was
$-10$ dB for 1 ms edges, $-24$ dB for 2 ms, $-40$ dB for 5 ms, and $-52$ dB for 10 ms.
\begin{parts}

	\begin{minipage}[t]{0.58\linewidth}
		\part A club station must keep its keying at least 30 dB below hard keying,
		1 kHz from the carrier. Which of the four edge times meet that, and which is the
		shortest that does?
		\begin{solution}[0.9in]
		The 5 ms and 10 ms edges meet it, at $-40$ and $-52$ dB. The 2 ms edge reaches only
		$-24$ dB. The shortest that meets the limit is 5 ms.
		\end{solution}
	\end{minipage}\hfill%
	\begin{minipage}[t]{0.37\linewidth}
		\raggedleft \vspace{12pt}
		\ansbox[1.1in]{5 ms}
	\end{minipage}

	\begin{minipage}[t]{0.58\linewidth}
		\part An operator wants 10 ms edges at 60 WPM. Find the dit length, and how much
		of each dit stays at full power. Each edge takes the full edge time, half of it
		before the nominal edge and half after.
		\begin{solution}[1.0in]
		\eq{ t_\text{dit} &= \dfrac{1.2\ \text{s}}{60} = 20\ \text{ms} \\
		20 - 10 &= 10\ \text{ms} }
		Half of each dit goes to the rising and falling edges.
		\end{solution}
	\end{minipage}\hfill%
	\begin{minipage}[t]{0.37\linewidth}
		\raggedleft \vspace{12pt}
		dit: \ansbox[1.0in]{20 ms}

		\vspace{4pt}
		full: \ansbox[1.0in]{10 ms}
	\end{minipage}

	\part In words, explain which receivers a hard-keyed transmitter interferes with, and why, and why the
	operator in part (b) should choose a shorter edge.
		\begin{solution}[1.2in]
		Hard keying spreads energy kilohertz away from the carrier, so a receiver tuned to
		another station nearby hears a click at every edge of the hard-keyed signal. The
		station itself decodes the same either way. At 60 WPM a 10 ms edge leaves only
		10 ms of each 20 ms dit at full power, so the dits lose their flat tops and a
		receiver in noise has less of each mark to detect. A 5 ms edge keeps
		15 ms of each dit at full power and still keeps the clicks well down.
		\end{solution}

\end{parts}

\newpage
%%%%%%%%%%%%%%%%%%%%%%%%%%%%%%%%%%%%%%%%%
\question \textbf{The product detector and the CW filter.} The L12 receiver moves the
tuned frequency to 0 Hz, multiplies by a 700 Hz BFO, and keeps the real part. Its CW
filter is a band-pass centered on 700 Hz. The noise has amplitude $\sigma$, spread over
32 kHz, and the keyed carrier has power 1, so the SNR in a filter of width $B$ is
\[ \text{SNR} = \frac{32{,}000}{\sigma^2 B}. \]
\begin{parts}

	\begin{minipage}[t]{0.58\linewidth}
		\part A cadet tunes the receiver 120 Hz above the carrier. What pitch does the tone
		have, and does it pass a 150 Hz CW filter? A 1 kHz filter?
		\begin{solution}[1.2in]
		Tuned 120 Hz high, the carrier lands at $-120$ Hz, so the tone is
		$700 - 120 = 580$ Hz. The 150 Hz filter passes 625 to 775 Hz, so the tone falls
		outside it and the decoder loses the signal. The 1 kHz filter passes 200 to
		1200 Hz, so the tone gets through.
		\end{solution}
	\end{minipage}\hfill%
	\begin{minipage}[t]{0.37\linewidth}
		\raggedleft \vspace{12pt}
		tone: \ansbox[1.0in]{580 Hz}

		\vspace{4pt}
		150 Hz: \ansbox[1.0in]{no}

		\vspace{4pt}
		1 kHz: \ansbox[1.0in]{yes}
	\end{minipage}

	\begin{minipage}[t]{0.58\linewidth}
		\part At $\sigma = 8$, find the SNR in the 200 Hz filter and in a 50 Hz filter.
		\begin{solution}[1.2in]
		\eq{ \text{SNR}_{200} &= \dfrac{32{,}000}{64 \times 200} = 2.5 = 4.0\db \\
		\text{SNR}_{50} &= \dfrac{32{,}000}{64 \times 50} = 10 = 10.0\db }
		\end{solution}
	\end{minipage}\hfill%
	\begin{minipage}[t]{0.37\linewidth}
		\raggedleft \vspace{12pt}
		200 Hz: \ansbox[1.0in]{4.0 dB}

		\vspace{4pt}
		50 Hz: \ansbox[1.0in]{10.0 dB}
	\end{minipage}

	\part Lesson 12's 200 Hz filter decoded without errors down to 10 dB. Explain in
	words why the 50 Hz filter in part (b) is the better choice in this noise, and what
	limits its use.
		\begin{solution}[1.2in]
		The signal is only about 100 Hz wide, so most of what a wider filter passes is
		noise. Cutting the width by four cuts the noise by four, 6 dB, which brings the SNR
		to the 10 dB where the 200 Hz filter decoded cleanly. Its limit is tuning: the tone
		must stay inside the narrow filter, so the receiver must tune within 25 Hz of the
		carrier, and a transmitter that drifts by more than that leaves the filter.
		\end{solution}

\end{parts}

\newpage
%%%%%%%%%%%%%%%%%%%%%%%%%%%%%%%%%%%%%%%%%
\question \textbf{Setting the threshold.} In the L12 receiver the Threshold block
switches the key on when the envelope rises above $1.2 \times$ \texttt{thr} and off when
it falls below $0.8 \times$ \texttt{thr}. On a noisy link, the envelope plot shows a
space level of about 0.30 and a mark level of about 0.70.
\begin{parts}

	\begin{minipage}[t]{0.58\linewidth}
		\part Choose \texttt{thr}, and find the two levels at which the key switches.
		\begin{solution}[1.0in]
		Halfway between the levels:
		\eq{ \texttt{thr} &= \dfrac{0.30 + 0.70}{2} = 0.50 \\
		\text{on} &= 1.2 \times 0.50 = 0.60, \quad \text{off} = 0.8 \times 0.50 = 0.40 }
		\end{solution}
	\end{minipage}\hfill%
	\begin{minipage}[t]{0.37\linewidth}
		\raggedleft \vspace{12pt}
		thr: \ansbox[1.0in]{0.50}

		\vspace{4pt}
		on / off: \ansbox[1.0in]{0.60 / 0.40}
	\end{minipage}

	\part A teammate leaves \texttt{thr} at 0.3, the value that worked with little noise.
	Explain what the key and the decoded text do.
		\begin{solution}[1.0in]
		The key switches on at 0.36, barely above the 0.30 space level, so noise peaks
		during the spaces turn it on. The decoder sees short extra marks, prints extra Es
		and Ts, and splits or corrupts letters, so most of the text comes out wrong.
		\end{solution}

	\part Another teammate sets \texttt{thr} to 0.8. Explain what happens, and why
	the hysteresis does not help.
		\begin{solution}[1.0in]
		The key now switches on only above 0.96, well above the 0.70 mark level, so only
		the strongest noise peaks during marks turn it on. The key misses most marks, and the
		decoder prints few letters, mostly short ones. Hysteresis only stops chatter near the
		threshold; it cannot fix a threshold that sits outside the range between the
		two levels.
		\end{solution}

\end{parts}

\vspace{1.5em}
\noindent\textbf{Documentation:} \rule[-2pt]{0.6\linewidth}{0.4pt}

\end{questions}
